\documentclass{article}
\usepackage{graphicx} % Required for inserting images

\title{ORIE 5370}

\begin{document}

\maketitle

\section{Extending Linear IPO to Nonlinear IPO}

\subsection{Motivation}

The original integrated prediction and optimization (IPO) framework considers a linear prediction model for expected asset returns and optimizes the prediction parameters directly with respect to downstream mean-variance portfolio performance. This is attractive because it aligns the training objective with the final investment decision rather than with a standalone prediction loss.

However, linear predictors may be too restrictive in practice. Asset returns often depend on features in a nonlinear manner, through interactions, higher-order effects, or more flexible latent representations. Therefore, a natural extension is to replace the linear return predictor with a nonlinear function while preserving the decision-focused IPO objective. This leads to a nonlinear IPO formulation.

There are two conceptually distinct routes to nonlinear IPO:
\begin{enumerate}
    \item \textbf{Nonlinear in inputs but linear in parameters:} use a nonlinear feature map $\phi(x)$ and define the predictor as a linear model in the transformed features.
    \item \textbf{Fully nonlinear in parameters:} use a neural network, kernel model, or another differentiable nonlinear predictor $f(x;\theta)$.
\end{enumerate}

The first route preserves much of the analytical structure of the original IPO formulation, while the second generally requires gradient-based optimization through the downstream portfolio problem.

\subsection{General Nonlinear IPO Formulation}

Let $x \in \mathbb{R}^{d_x}$ denote the input features and let $y \in \mathbb{R}^{d_y}$ denote the realized asset returns. We consider a general predictor
\[
\hat y = f(x;\theta),
\]
where $\theta$ denotes the model parameters.

For each observation $i=1,\dots,m$, the downstream portfolio decision is obtained by solving a mean-variance optimization problem using the predicted returns $\hat y_i = f(x_i;\theta)$. The nonlinear IPO problem is
\begin{equation}
\min_{\theta} \; \frac{1}{m}\sum_{i=1}^{m} c\!\left(z_i^*(f(x_i;\theta)),\, y_i\right)
\label{eq:nonlinear_ipo_outer}
\end{equation}
subject to
\begin{equation}
z_i^*(\hat y_i)
=
\arg\min_{z \in S}
\left(
-z^\top \hat y_i
+ \frac{\delta}{2} z^\top \hat V_i z
\right),
\label{eq:nonlinear_ipo_inner}
\end{equation}
where $\delta > 0$ is the risk-aversion parameter, $\hat V_i$ is the estimated covariance matrix, and $S$ denotes the feasible set of portfolio weights.

The realized downstream cost is
\begin{equation}
c(z,y)
=
-z^\top y + \frac{\delta}{2} z^\top V z,
\label{eq:realized_cost}
\end{equation}
where $V$ is the true covariance matrix (or an out-of-sample proxy).

\subsection{Route I: Nonlinear Features with Parameter-Linear Prediction}

A natural extension is to introduce a nonlinear feature map
\[
\phi(x) \in \mathbb{R}^{p},
\]
and define the prediction model as
\begin{equation}
\hat y = W \phi(x) + b.
\label{eq:feature_model}
\end{equation}
Examples include polynomial expansions, interaction terms, spline bases, radial basis features, or other engineered nonlinear transformations.

To simplify notation, augment the feature vector with a constant:
\[
\tilde \phi(x)
=
\begin{bmatrix}
\phi(x) \\
1
\end{bmatrix},
\qquad
\hat y = \Theta^\top \tilde \phi(x),
\]
where $\Theta$ stacks $W$ and $b$.

Although the predictor is nonlinear in the original inputs $x$, it remains linear in the parameters $\Theta$. Therefore, the analytical structure of the original linear IPO framework can be preserved after replacing $x$ with $\phi(x)$.

\subsection{Unconstrained Mean-Variance Problem}

We first consider the unconstrained case $S=\mathbb{R}^{d_y}$. The inner problem becomes
\[
z_i^*(\hat y_i)
=
\arg\min_{z}
\left(
-z^\top \hat y_i + \frac{\delta}{2} z^\top \hat V_i z
\right).
\]
The first-order condition is
\[
-\hat y_i + \delta \hat V_i z = 0,
\]
which implies
\begin{equation}
z_i^*(\hat y_i)
=
\frac{1}{\delta}\hat V_i^{-1}\hat y_i.
\label{eq:unconstrained_optimal_z}
\end{equation}

Substituting the parameter-linear nonlinear predictor $\hat y_i=\Theta^\top \tilde\phi_i$ into \eqref{eq:unconstrained_optimal_z}, we obtain
\[
z_i^*(\Theta)
=
\frac{1}{\delta}\hat V_i^{-1}\Theta^\top \tilde\phi_i.
\]

The realized cost for sample $i$ is
\[
\ell_i(\Theta)
=
-z_i^{*\top} y_i
+\frac{\delta}{2} z_i^{*\top} V_i z_i^*.
\]
Substituting $z_i^*(\Theta)$ gives
\begin{align}
\ell_i(\Theta)
&=
-\frac{1}{\delta}\tilde\phi_i^\top \Theta \hat V_i^{-1} y_i
+
\frac{1}{2\delta}
\tilde\phi_i^\top
\Theta \hat V_i^{-1} V_i \hat V_i^{-1}
\Theta^\top \tilde\phi_i.
\label{eq:loss_theta_unconstrained}
\end{align}

Therefore, the full nonlinear-feature IPO objective is
\begin{equation}
L(\Theta)
=
\frac{1}{m}\sum_{i=1}^m \ell_i(\Theta),
\label{eq:full_loss_unconstrained}
\end{equation}
which is a quadratic function of $\Theta$. Hence, this extension remains analytically tractable and can be solved as a convex quadratic optimization problem, or equivalently as a generalized normal equation system.

This shows that \emph{nonlinearity in the input space does not destroy the analytical IPO structure as long as the predictor remains linear in its parameters}.

\subsection{Equality-Constrained Mean-Variance Problem}

Next, suppose the feasible set is defined by linear equality constraints,
\[
S = \{ z : Az = b \}.
\]
Using a null-space parameterization, any feasible portfolio can be written as
\[
z = Fq + z_0,
\]
where the columns of $F$ span the null space of $A$, i.e.\ $AF=0$, and $z_0$ is any feasible point satisfying $Az_0=b$.

Substituting into the inner problem yields an optimization problem in $q$. The optimal decision can be expressed as
\begin{equation}
z_i^*(\hat y_i)
=
\frac{1}{\delta}
F(F^\top \hat V_i F)^{-1}F^\top \hat y_i
+
\left(
I - F(F^\top \hat V_i F)^{-1}F^\top \hat V_i
\right)z_0.
\label{eq:equality_solution}
\end{equation}

Define
\[
B_i
=
\frac{1}{\delta}F(F^\top \hat V_i F)^{-1}F^\top,
\qquad
a_i
=
\left(
I - F(F^\top \hat V_i F)^{-1}F^\top \hat V_i
\right)z_0.
\]
Then
\begin{equation}
z_i^*(\hat y_i)=B_i \hat y_i + a_i.
\label{eq:affine_decision_rule}
\end{equation}

Substituting $\hat y_i=\Theta^\top \tilde\phi_i$ into \eqref{eq:affine_decision_rule}, the decision becomes affine in $\Theta$. The sample loss is
\[
\ell_i(\Theta)
=
-(B_i\Theta^\top \tilde\phi_i + a_i)^\top y_i
+\frac{\delta}{2}(B_i\Theta^\top \tilde\phi_i + a_i)^\top
V_i
(B_i\Theta^\top \tilde\phi_i + a_i).
\]
Again, this is quadratic in $\Theta$, so the nonlinear-feature IPO formulation remains analytically solvable in the equality-constrained case.


\end{document}
